% -------------------------------------------------------
\section{Conclusion}
\label{sec:rs-conclusion}

The Recursive Spawning Protocol makes autonomous drift architecturally impossible — not by prescribing agent behavior, not by enforcing policy conventions, but by establishing immutable parent pointers as a first-class architectural foundation of multi-agent accountability. This is not a refinement of existing oversight approaches. It is a structural redefinition of what a delegation chain \emph{is}: from a mutable reference subject to retroactive rewriting, to a runtime artifact that is immutable from the moment of provisioning and auditable at any future point without trusting any machine actor in the chain.

\medskip
\noindent\textbf{Immutability as architectural foundation.} Every agent $n$ is provisioned with $\alpha(n)$ set once, at provisioning, by the Fact Layer write protocol. No actor — under any operational circumstance, with any level of privilege, for any stated justification — may modify $\alpha(n)$ post-provisioning. The parent pointer is not access-controlled; it is structurally immutable by protocol definition. This single constraint eliminates all four drift failure modes by construction. Chain opacity cannot arise because the delegation chain is not a mutable reference — it is a first-class runtime artifact established by the Fact Layer write protocol and hash-chained into the tamper-evident forensic ledger. Scope creep drift cannot accumulate because the Fact Layer pre-commit hook enforces $R(n_{i+1}) \subseteq R(n_i)$ as a structural invariant, not a policy convention. Re-parenting drift is architecturally impossible because retroactive pointer modification is structurally forbidden. Ghost authority cannot propagate because every spawn requires a Purpose Layer authorization record in the forensic ledger before the child node is provisioned.

\medskip
\noindent\textbf{Three-layer structural substrate.} The \bxthree three-layer architecture provides the structural substrate that makes the immutable parent pointer effective — each layer serving a role exclusive to it. The Purpose Layer defines spawn authorization policy and establishes the human principal as the exclusive terminus of every delegation chain. The Bounds Engine enforces depth constraints as hard architectural boundaries, preventing unbounded spawn growth that would exceed any human's supervisory capacity. The Fact Layer enforces the immutable parent pointer via its write protocol, maintains the tamper-evident forensic ledger, and operates the pre-commit hook that converts scope refinement from policy documentation into structural enforcement. The mapping is minimal-sufficient: no layer is redundant, and no proper subset of layers achieves the same guarantees.

\medskip
\noindent\textbf{Correctness properties.} The three correctness properties are jointly operative and mutually reinforcing. Drift prevention (Theorem 1) ensures each child's operating scope is a descendant refinement of its parent's, traceable to the human anchor's intent through the forensic ledger — making scope divergence architecturally impossible regardless of chain depth. Chain integrity (Theorem 2) ensures the chain's topology is stable and its parent pointers unmodifiable after provisioning, with any tampering detectable by hash-chain verification. Termination (Theorem 3) ensures the chain halts within a human-defined depth bound and can sustain no infinite spawn-escalate cycle. Together they constitute the RSP's full operational guarantee: a spawn chain is an accountable, bounded, and auditable refinement process that terminates in human judgment.

\medskip
\noindent\textbf{Known costs, identified mitigations.} The immutable parent pointer prevents legitimate chain reconfiguration; a successor promotion mechanism is the identified resolution (Section~\ref{sec:rs-future-work}). The pre-commit hook and hash-chained forensic ledger impose runtime overhead quantified by Equation~\ref{eq:spawn-overhead}, mitigated by checkpointing and parallel write optimizations. The protocol's scope is intra-chain; inter-chain accountability in multi-tenant deployments requires co-deployment with Spatial Firewall and Sandbox Gate and a unified correctness statement that RSP alone does not provide. The compromised Purpose Layer risk is load-bearing; it is addressed by Purpose Layer governance and provisioning-time integrity monitoring, not by RSP modification. None of these limitations undermines the core correctness properties under the stated assumptions.

\medskip
\noindent\textbf{Contribution to the accountability landscape.} RSP is not a monitoring tool added to an existing runtime; it is a structural redefinition of the delegation chain itself. It is not a depth limit with better configuration; it is a proof (Section~\ref{sec:drift-depth-limits}) that depth limits alone cannot prevent drift, and that the conjunction of immutable parent pointers, scope refinement as a structural invariant, and human-exclusive chain termination is both necessary and sufficient for drift prevention. It is not organizational policy; it is an architectural constraint that holds regardless of the goodness, reliability, or intentions of any machine actor in the chain.

\medskip
\noindent\textbf{The upstream guarantee, realized.} The upstream \bxthree accountability guarantee — that systems built on the \bxthree Framework fail upward into human consciousness, never downward into autonomous chaos — is given its most complete expression in the Recursive Spawning Protocol. When an agent in an RSP-governed system fails, the failure surfaces to a human principal through a complete, immutable, tamper-evident accountability chain. When an RSP-governed system encounters a condition it cannot resolve autonomously, it compulsory-escalates. There is no mechanism by which the system can proceed without human judgment, because the protocol is structured to make human judgment the non-negotiable terminus of every delegation chain.

The Recursive Spawning Protocol is the second named pillar of the \bxthree Framework. It is available for integration into any agentic system that requires the structural guarantee that every agent's authority traces to a human principal — not by policy assertion, not by forensic reconstruction, but by architectural construction from the moment of provisioning.
